La figura mostra l'estructura d'un accelerador lineal. Els cilindres metàl·lics tenen càrregues elèctriques alternades positives i negatives. A l'interior dels cilindres el camp elèctric és negligible. Les partícules carregades són accelerades per un camp elèctric aproximadament uniforme als espais entre els tubs ($A$, $B$, $C$\ldots). La diferència de potencial entre cilindres és de 250~kV.

\figura[0.5]{fig1.png}

\begin{apartat}{1,25}
Si volem accelerar un electró que es mou del primer al segon cilindre, quins signes hauran de tenir les càrregues acumulades al primer i al segon cilindre? Justifiqueu la resposta. Dibuixeu les línies de camp elèctric a l'espai $A$ entre cilindres. Si volem obtenir un camp de $8,00\times 10^{6}\ \mathrm{N/C}$, calculeu la distància que separa els dos primers cilindres.
\end{apartat}

\begin{apartat}{1,25}
Per a mantenir el sentit de l'acceleració, les polaritats dels cilindres s'inverteixen cada vegada que l'electró entra en el cilindre següent. Quants espais entre cilindres hauria de tenir l'accelerador si volem que l'electró surti amb una energia d'1,0~MeV? Sense tenir en compte la correcció relativista, quina seria la velocitat dels electrons? Comenteu el resultat obtingut.
\end{apartat}

\dades{%
  $1\ \mathrm{eV} = 1,602\times 10^{-19}\ \mathrm{J}$.\\
  $|e| = 1,602\times 10^{-19}\ \mathrm{C}$.\\
  $m_\mathrm{e} = 9,11\times 10^{-31}\ \mathrm{kg}$.\\
  $c = 3,00\times 10^{8}\ \mathrm{m\,s^{-1}}$.}
